\notechapter{N-20260313}{Multivariable Chain Rule, Total Differentials, and Implicit Differentiation}

\section{The chain rule is composition of linear maps}

The chain rule is not a collection of unrelated formulas.  It says that first-order approximations compose.

\begin{theorem}[Chain rule]
Let \(F:\mathbb R^n\to\mathbb R^m\) be differentiable at \(x\), and let \(G:\mathbb R^m\to\mathbb R^k\) be differentiable at \(F(x)\).  Then \(G\circ F\) is differentiable at \(x\), and
\[
  D(G\circ F)(x)=DG(F(x))\,DF(x).
\]
\end{theorem}

\begin{proof}[Proof skeleton]
Write
\[
  F(x+h)=F(x)+DF(x)h+r(h),\qquad \norm{r(h)}=o(\norm h).
\]
Then apply the differentiability expansion for \(G\) at \(F(x)\):
\[
  G(F(x)+u)=G(F(x))+DG(F(x))u+s(u),
  \qquad \norm{s(u)}=o(\norm u).
\]
Substitute \(u=DF(x)h+r(h)\).  Since \(u=O(\norm h)\), the remainder is still \(o(\norm h)\).  The linear part is \(DG(F(x))DF(x)h\).
\end{proof}

This theorem is the source of all ordinary partial-derivative chain rules.

\section{One input variable, two intermediate variables}

Suppose
\[
  z=f(u,v),\qquad u=u(x),\qquad v=v(x).
\]
Then
\[
  \frac{dz}{dx}=f_u\frac{du}{dx}+f_v\frac{dv}{dx}.
\]
In matrix form, this is
\[
  \frac{dz}{dx}
  =
  \begin{pmatrix}f_u&f_v\end{pmatrix}
  \begin{pmatrix}u_x\\ v_x\end{pmatrix}.
\]
The row vector is \(Df\), the column vector is the derivative of \((u,v)\), and their product is the derivative of the composite.

\begin{example}
Let \(z=u^2+uv\), \(u=x^2\), and \(v=\sin x\).  Then
\[
  z_u=2u+v,\qquad z_v=u,\qquad u'=2x,\qquad v'=\cos x.
\]
Hence
\[
  \frac{dz}{dx}=(2u+v)2x+u\cos x.
\]
Substituting \(u=x^2\), \(v=\sin x\) gives
\[
  \frac{dz}{dx}=2x(2x^2+\sin x)+x^2\cos x.
\]
\end{example}

\section{Two input variables}

Suppose
\[
  z=f(u,v),\qquad u=u(x,y),\qquad v=v(x,y).
\]
Then
\[
  \begin{pmatrix}z_x&z_y\end{pmatrix}
  =
  \begin{pmatrix}f_u&f_v\end{pmatrix}
  \begin{pmatrix}
    u_x&u_y\\
    v_x&v_y
  \end{pmatrix}.
\]
Equivalently,
\[
  z_x=f_u u_x+f_v v_x,\qquad
  z_y=f_u u_y+f_v v_y.
\]

\begin{example}
Let \(z=e^u\cos v\), where \(u=x^2+y^2\) and \(v=x-y\).  Then
\[
  f_u=e^u\cos v,\qquad f_v=-e^u\sin v,
\]
and
\[
  \begin{pmatrix}
    u_x&u_y\\ v_x&v_y
  \end{pmatrix}
  =
  \begin{pmatrix}
    2x&2y\\ 1&-1
  \end{pmatrix}.
\]
Thus
\[
  z_x=e^u\cos v\cdot2x-e^u\sin v,
  \qquad
  z_y=e^u\cos v\cdot2y+e^u\sin v.
\]
\end{example}

\section{Total differentials}

The notation
\[
  dz=f_x\,dx+f_y\,dy
\]
is best read as a compact way to write a linear approximation.  If \((x,y)\) changes by \((dx,dy)\), then the first-order change in \(z=f(x,y)\) is
\[
  dz\approx f_x\,dx+f_y\,dy.
\]
For a vector-valued map \(F\), the same statement is
\[
  dF=J_F\,dx.
\]
The differential notation is safe when one remembers that it represents a linear map, not ordinary fraction arithmetic.

\section{Implicit differentiation in one equation}

Suppose a relation
\[
  F(x,y)=0
\]
defines \(y\) locally as a differentiable function of \(x\).  Differentiating the identity \(F(x,y(x))=0\) gives
\[
  F_x+F_y\frac{dy}{dx}=0.
\]
If \(F_y\ne0\), then
\[
  \frac{dy}{dx}=-\frac{F_x}{F_y}.
\]

\begin{example}
For the circle \(F(x,y)=x^2+y^2-1=0\), where \(y\ne0\),
\[
  \frac{dy}{dx}=-\frac{2x}{2y}=-\frac{x}{y}.
\]
This is the slope of the tangent line to the circle.
\end{example}

\section{Implicit differentiation for systems}

The same idea becomes matrix algebra.  Suppose
\[
  F(x,u)=0,\qquad
  F:\mathbb R^p\times\mathbb R^q\to\mathbb R^q,
\]
and \(u\) is locally a function of \(x\).  Differentiating gives
\[
  F_x\,dx+F_u\,du=0.
\]
If \(F_u\) is invertible, then
\[
  du=-F_u^{-1}F_x\,dx,
\]
so
\[
  Du(x)=-F_u^{-1}F_x.
\]

\begin{theorem}[Implicit function theorem, working form]
Let \(F:\mathbb R^p\times\mathbb R^q\to\mathbb R^q\) be \(C^1\), and suppose
\[
  F(a,b)=0,\qquad \det F_u(a,b)\ne0.
\]
Then near \((a,b)\), the equation \(F(x,u)=0\) determines \(u\) as a \(C^1\) function of \(x\), and
\[
  Du(a)=-F_u(a,b)^{-1}F_x(a,b).
\]
\end{theorem}

The theorem says when the formal linearized computation is legitimate.  The condition \(\det F_u\ne0\) means the constraint can be solved in the \(u\)-directions.

\section{A shape table}

Keeping track of shapes prevents many mistakes:
\[
\begin{array}{c|c|c}
\text{map} & \text{derivative shape} & \text{meaning}\\
\hline
f:\mathbb R^n\to\mathbb R & 1\times n & \text{row gradient}\\
F:\mathbb R^n\to\mathbb R^m & m\times n & \text{Jacobian}\\
G\circ F & k\times n & (k\times m)(m\times n)\\
F(x,u)=0 & q\times p,\ q\times q & F_x,F_u
\end{array}
\]
The formulas work because the linear maps have compatible domains and codomains.

\section{Singular linearization marks the boundary of implicit solvability}

The invertibility condition in the implicit function theorem cannot simply be omitted.  Consider
\[
  F(x,y)=y^2-x.
\]
At the origin,
\[
  F(0,0)=0,
  \qquad
  F_y(0,0)=0.
\]
The zero set is
\[
  x=y^2,
\]
or, for \(x\ge0\), the two branches
\[
  y=\sqrt{x},
  \qquad
  y=-\sqrt{x}.
\]
No neighborhood of the origin in the zero set is the graph of a single function \(y=u(x)\), and the two branch derivatives become unbounded as \(x\downarrow0\).  The formal formula
\[
  u'(x)=-\frac{F_x}{F_y}=\frac1{2y}
\]
therefore signals its own breakdown at \(y=0\).

The same zero set can be solved smoothly in the other direction:
\[
  x=y^2,
\]
because
\[
  F_x(0,0)=-1\ne0.
\]
Implicit solvability is directional.  The Jacobian block that is invertible determines which variables may be expressed as functions of the others.
